\vspace*{4cm}
\begin{center}
    \begin{Arabic}
        {\Huge\bfseries الملخص}
    \end{Arabic}
\end{center}
\vskip 1cm

\phantomsection
\addcontentsline{toc}{chapter}{\texorpdfstring{\textarabic{الملخص}}{الملخص}}

\begin{Arabic}
أدت التحولات الرقمية في قطاع الصحة إلى ظهور الأجهزة الطبية المتصلة، والسجلات الصحية الإلكترونية،
ومنصات المراقبة عن بعد، وأنظمة دعم القرار المعتمدة على الذكاء الاصطناعي. ورغم أن هذه التقنيات
تنتج كميات كبيرة من البيانات المتنوعة، فإن العديد من أنظمة الصحة الإلكترونية ما تزال مجزأة، وثابتة،
ومحدودة التخصيص. وتقدم تقنية التوأم الرقمي مقاربة واعدة لمعالجة هذه الحدود، من خلال إنشاء
تمثيلات افتراضية ديناميكية للكيانات الفيزيائية أو العمليات أو المرضى، يتم تحديثها باستمرار عبر تبادل
البيانات، وتدعيمها بالمحاكاة والتحليل وخدمات دعم القرار.

يعرض هذا البحث دراسة لحالة الفن حول توظيف التوائم الرقمية في مجال الصحة الإلكترونية. ويتناول
أولا الأسس العامة للصحة الإلكترونية، وتطور مفهوم التوأم الرقمي، والتمييز بين النموذج الرقمي، والظل
الرقمي، والتوأم الرقمي المتكامل. ثم يحلل تطبيقات التوائم الرقمية الموجهة للرعاية الصحية، ومعمارياتها،
والتقنيات الممكنة لها، وتحديات تنفيذها، مع توسيع النقاش ليشمل مجالات تطبيقية مثل الطب الشخصي،
والمراقبة عن بعد، وتحسين سير العمل داخل المستشفيات، والصحة العامة، والتوائم الرقمية القلبية الوعائية.
كما يولي البحث اهتماما خاصا لقضايا قابلية التشغيل البيني، وحوكمة البيانات، والخصوصية، والأمن،
وقابلية التفسير، والتحقق، والاعتبارات الأخلاقية. وفي الأخير،
يقترح البحث معمارية مفاهيمية متعددة الطبقات لنظام صحة إلكترونية قائم على التوأم الرقمي، بهدف دعم
المراقبة المستمرة، والتحليل الشخصي، ومحاكاة السيناريوهات، والمساعدة على اتخاذ القرار في سياق صحي
عام.

\textbf{الكلمات المفتاحية:} التوأم الرقمي، الصحة الإلكترونية، الصحة الرقمية، الرعاية الصحية،
إنترنت الأشياء، إنترنت الأشياء الطبية، الذكاء الاصطناعي، المحاكاة، قابلية التشغيل البيني، دعم القرار
السريري، التوأم الرقمي القلبي الوعائي.
\end{Arabic}

\clearpage
